\newpage
\appendix

\section{Identifiability of GRL-SNAM}
\label{app:proof:and:lemma}
\subsection{From optimal control to reduced Hamiltonian dynamics }
\label{app:ocp_to_ham}
\begin{lemma}[Forward: autonomous pure state constraints induce a (reduced) Hamiltonian flow]
\label{thm:forward_state_constrained_pmp_reduced_ham}
Fix a scenario $\mathcal{E}$ and consider the autonomous, control-affine system on an open set
$\mathcal{Q}\subset\mathbb{R}^n$,
\begin{equation}
\dot q(t)=f(q(t)) + A(q(t))u(t),\qquad q(0)=q_0,\qquad q(T)=q_T,
\label{eq:ocp_dyn}
\end{equation}
with pure state inequality constraints
\begin{equation}
d_i(q(t))\ge 0,\qquad i=1,\dots,m,\qquad t\in[0,T],
\label{eq:ocp_state_constraints}
\end{equation}
and objective
\begin{equation}
\min_{u(\cdot)\in\mathcal{U}} \ J(q, u)
:=\int_0^T L(q(t),u(t))\,dt,
\label{eq:ocp_cost}
\end{equation}
where the admissible control set $\mathcal{U}$ consists of measurable $u(\cdot)$ with
$u(t)\in U\subset\mathbb{R}^d$ a.e.

Assume:
\begin{enumerate}
\item $f$ and $A$ are locally Lipschitz on $\mathcal{Q}$.
\item $U$ is nonempty, compact, and convex; $L(\cdot,\cdot)$ is continuous and $C^1$ in $(q,u)$.
\item (Existence) There exists an optimal pair $(q^\star(\cdot),u^\star(\cdot))$ satisfying
\cref{eq:ocp_dyn} and \cref{eq:ocp_state_constraints}.
\item (Regularity / qualification) Each $d_i\in C^2(\mathcal{Q})$ and along boundary arcs the active constraint
gradients satisfy a standard constraint qualification (e.g.\ linear independence of
$\{\nabla d_i(q^\star(t))\}_{i\in\mathcal{A}(t)}$ and a feasible-direction condition), so that the
state-constrained maximum principle applies.
\end{enumerate}

Define the \emph{control Hamiltonian} as \textit{Fenchel Coupling}
\begin{equation}
\mathcal{H}_c(q,p,u):=\langle f(q)+A(q)u, p\rangle-L(q,u).
\label{eq:def_control_ham}
\end{equation}
Then there exist:
(i) a costate $p:[0,T]\to\mathbb{R}^n$ of bounded variation(\textbf{BV}, right-continuous with left limits),
and (ii) nonnegative finite Radon measures $\{\mu_i\}_{i=1}^m$ on $[0,T]$,
such that the following hold:

\smallskip
\noindent\textbf{(i) Maximum condition (a.e.).}
For almost every $t\in[0,T]$,
\begin{equation}
u^\star(t)\in \arg\max_{u\in U}\ \mathcal{H}_c\big(q^\star(t),p(t),u\big).
\label{eq:max_condition}
\end{equation}
\noindent\textbf{(ii) State equation.}
For a.e.\ $t$,
\begin{equation}
\dot q^\star(t)=\nabla_p \mathcal{H}_c\big(q^\star(t),p(t),u^\star(t)\big)
= f(q^\star(t))+A(q^\star(t))u^\star(t).
\label{eq:state_eq_from_ham}
\end{equation}
\noindent\textbf{(iii) Adjoint equation in measure form.}
As an equality of $\mathbb{R}^n$-valued measures on $[0,T]$,
\begin{equation}
dp(t)= -\nabla_q \mathcal{H}_c\big(q^\star(t),p(t),u^\star(t)\big)\,dt
\;-\;\sum_{i=1}^m \nabla d_i\big(q^\star(t)\big)\,\mu_i(dt).
\label{eq:adjoint_measure}
\end{equation}
Equivalently, for all $0\le t_0<t_1\le T$,
\begin{equation}
p(t_1^+)-p(t_0^-)
= -\int_{t_0}^{t_1}\nabla_q \mathcal{H}_c\big(q^\star(t),p(t),u^\star(t)\big)\,dt
\;-\;\sum_{i=1}^m \int_{(t_0,t_1]} \nabla d_i\big(q^\star(t)\big)\,\mu_i(dt).
\label{eq:adjoint_integral}
\end{equation}
\noindent\textbf{(iv) Complementarity / support.}
Each $\mu_i$ is supported on the active set:
\begin{equation}
\mathrm{supp}(\mu_i)\subset \{t\in[0,T]: d_i(q^\star(t))=0\},
\qquad\text{i.f.f.}\qquad
\int_{[0,T]} d_i(q^\star(t))\,\mu_i(dt)=0.
\label{eq:complementarity_measure}
\end{equation}
\noindent\textbf{(iv) Reduced Hamiltonian.} Assume in addition that the running cost decomposes as
\begin{equation}
L(q,u)=\varphi(u)+V(q),
\qquad\text{with $\varphi$ proper, closed, strictly convex, and $C^1$, and $V\in C^1(\mathcal Q)$.}
\label{eq:cost_split_B_statecon}
\end{equation}
Then the maximizer $u^\star(q,p)$ in \cref{eq:max_condition} is unique (up to the compact control constraint $U$),
and the reduced PMP Hamiltonian depends only on $(q,p)$:
\begin{equation}
\label{eq:def_reduced_ham_pmp_statecon}
H_{\mathrm{PMP}}(q,p)
:=\sup_{u\in U}\mathcal H_c(q,p,u)
=\sup_{u\in U}\big\{p^\top A(q)u-\varphi(u)\big\}+p^\top f(q)-V(q).
\end{equation}
Hence \cref{eq:state_eq_from_ham} and \cref{eq:adjoint_measure} can be restated without explicit dependency on $u^\star$.
\end{lemma}

\begin{proof}
We present the derivation in steps and refer readers Theorem 4.2 in \cite{hartl1995survey} for further discussions. The proof contains major steps including reducing the general state-constrained PMP statement
to the present autonomous, pure-state-constraint setting and then eliminating $u$ via conjugacy.

\smallskip
\noindent\textbf{Step 1 (Apply the state-constrained maximum principle).}
Under assumptions (1)--(4), the state-constrained maximum principle provides a costate of bounded variation and
a bounded-variation multiplier associated with the state constraints, together with an integral adjoint relation.
Specializing to \emph{pure} state constraints (no mixed constraints) yields an identity of the form
\cref{eq:adjoint_integral} except that the state-constraint multiplier appears as a BV function $\zeta_i(\cdot)$.
Writing its Stieltjes increment as $d\zeta_i(t)$ and using the BV decomposition, define the nonnegative measure
$\mu_i := -d\zeta_i$; then $\mu_i\ge 0$ and \cref{eq:adjoint_integral} follows.

\smallskip
\noindent\textbf{Step 2 (Maximum condition from needle variations).}
Fix a Lebesgue point $t$ for $(q^\star,u^\star)$ and consider a needle variation
$u_\varepsilon(\cdot)$ that equals $u^\star(\cdot)$ except on a small interval $(t,t+\varepsilon)$ where it is set
to a constant $\bar u\in U$.
Using (A1)--(A2), the first variation of $J$ exists and is linear in $\bar u-u^\star(t)$, while feasibility of
\cref{eq:ocp_state_constraints} is handled by the measure multipliers $\mu_i$.
The necessary condition that the first variation is nonnegative for all $\bar u\in U$ implies
\cref{eq:max_condition} for almost every $t$.

\smallskip
\noindent\textbf{Step 3 (Adjoint equation from state variations + integration by parts).}
Let $\delta q$ be an admissible first-order state variation induced by admissible control variations.
Write the first variation of the Lagrangian functional associated with \cref{eq:ocp_dyn} using a multiplier $p(\cdot)$,
and integrate by parts on each interval of continuity of $p$.
The terms involving the state constraints appear through Stieltjes integrals against $d\zeta_i$, i.e.,
through $\int \nabla c_i(q^\star(t))\,d\zeta_i(t)$; converting to $\mu_i=-d\zeta_i$ yields \cref{eq:adjoint_measure}.

\smallskip
\noindent\textbf{Step 4 (Complementarity/support).}
Because each $\zeta_i(\cdot)$ is constant on interior intervals where $c_i(q^\star(t))>0$,
its Stieltjes increment vanishes there; hence $\mu_i=-d\zeta_i$ only charges times when the constraint is active.
This is exactly \cref{eq:complementarity_measure}.

\smallskip
\noindent\textbf{Step 5 (Reduced Hamiltonian by Legendre--Fenchel conjugacy).}
~Define the \emph{reduced Hamiltonian} (pointwise supremum over $u$)
\begin{equation}
H_{\mathrm{PMP}}(q,p)
:=\sup_{u\in U}\Big\{p^\top A(q)u-\varphi(u)\Big\}+p^\top f(q)-V(q).
\end{equation}
By strict convexity of $\varphi$, the maximizer $u^\star(q,p)$ is unique and satisfies the first-order condition
$\nabla\varphi(u^\star)=A(q)^\top p$ (when $U=\mathbb{R}^d$; for compact $U$ use the argmax form).
Hence $u^\star(q,p)=\nabla\varphi^\ast(A(q)^\top p)$, and Danskin's theorem yields
\[
\begin{aligned}
\nabla_p H_{\mathrm{PMP}}(q,p)&=f(q)+A(q)u^\star(q,p), \\
\nabla_q H_{\mathrm{PMP}}(q,p)&=\big(\nabla_q f(q)\big)^\top p+\big(\nabla_q A(q)\,u^\star(q,p)\big)^\top p-\nabla V(q).    
\end{aligned}
\]
Substituting $u^\star(t)=u^\star(q^\star(t),p(t))$ into
\cref{eq:state_eq_from_ham}--\cref{eq:adjoint_measure} gives the \emph{reduced} constrained Hamiltonian flow
\begin{equation}
\begin{aligned}
\dot q^\star(t)&=\nabla_p H_{\mathrm{PMP}}\big(q^\star(t),p(t)\big), \\ 
dp(t)&= -\nabla_q H_{\mathrm{PMP}}\big(q^\star(t),p(t)\big)\,dt -\sum_{i=1}^m \nabla d_i(q^\star(t))\,\mu_i(dt),    
\end{aligned}
\end{equation}
which is precisely the claimed Hamiltonian outcome of the state-constrained optimal control problem.
\end{proof}

\begin{corollary}[Reduced Hamiltonian via a quadratic control penalty]
\label{cor:mechanical_reduced_hamiltonian}
Consider the control--affine dynamics $\dot q=f(q;\mathcal E)+A(q;\mathcal E)u$ and assume the { cost-to-go functional} splits as
$L(q,u;\mathcal E)=\varphi(u) + {V}(q;\mathcal E)$ with a quadratic effort penalty
\(
\varphi(u)=\tfrac12 u^\top \Phi u
\)
 for some $\Phi\succ 0$. Assume $f(\cdot;\mathcal E)\equiv 0$ and $A(q;\mathcal{E})\equiv A$.
Then the reduced Hamiltonian (i.e.\ $\sup_u\{p^\top(f+A u)-L\}$) equals
\[
H(q,p;\mathcal E)
=\tfrac12\,p^\top\!\big(A(q;\mathcal E)\Phi^{-1}A(q;\mathcal E)^\top\big)p -V(q;\mathcal E).
\]
In particular, if we define the (inverse) mass by
\[
M(q;\mathcal E)^{-1}:=A(q;\mathcal E)\Phi^{-1}A(q;\mathcal E)^\top\in\mathbb S_{++}^2,
\]
then \(H(\cdot,\cdot;\mathcal E)\in\mathscr H\), where
\[
\mathscr{H}_{PMP}:=\Big\{\,H(q,p;\mathcal{E})=\tfrac12\,p^\top M^{-1}p{-}V(q;\mathcal{E})
\ \Big|\ M\in \mathbb{S}_{++}^2 \Big\}.
\]
\end{corollary}

\begin{proof}
By definition of the PMP reduced Hamiltonian,
\[
H_{PMP}(q,p;\mathcal E)=\sup_{u}\Big\{p^\top(f(q;\mathcal E)+A(q;\mathcal E)u)-\varphi(u)- V(q;\mathcal E)\Big\}.
\]
Setting $f\equiv 0$ reduces the supremum to$\sup_u\{(A^\top p)^\top u-\tfrac12 u^\top \Phi u\}{-V}$.
The maximizer satisfies $\Phi u^\star=A^\top p$, hence $u^\star=\Phi^{-1}A^\top p$ and the supremum value is
$\tfrac12 p^\top A\Phi^{-1}A^\top p { - V}$. Since $\Phi\succ 0$ and $A(\cdot;\mathcal E)$ has full rank, the induced
matrix $A\Phi^{-1}A^\top$ is positive definite, so $M^{-1}\in\mathbb S_{++}^2$ and $H\in\mathscr H$.
\end{proof}
{
\paragraph{Barrier relaxation of state-constraint OCPs}\label{rem:barrier_relaxation}
In the implementation one replaces hard constraints $d_i(q)\ge 0$ by a smooth barrier/penalty term
$B_\mu(q):=\mu \sum_i b(d_i(q))$ and optimizes an \emph{unconstrained} problem with augmented running cost
$L_\mu(q,u,t)=L(q,u,t)+B_\mu(q)$.
This does not require introducing BV multipliers in the PMP dynamics; instead it modifies the potential (equivalently, modifies $\nabla_q H_{PMP}$) by adding $\nabla_q B_\mu$. From the inverse-design viewpoint, learning barrier weights corresponds to selecting a particular absolutely continuous approximation of the constraint measure $d\zeta$ (as a standard exact-penalty / interior-point interpretation). For our GRL-SNAM navigator, we provide additional statements (c.f. \cref{prop:barrier_to_state_constrained_complete} and \cref{cor:ipc_from_logbarrier}) to address the approximability of our choice of Hamiltonians.
}

\begin{proposition}[Barrier relaxation $\mu\downarrow 0$ yields a state-constrained limit]
\label{prop:barrier_to_state_constrained_complete}
Fix $\mathcal E$ and consider the deterministic state-constrained OCP
\begin{equation}
\begin{aligned}
\label{eq:OCP_hard}
\min_{u(\cdot)}&\ J(u):=\int_0^T L(q(t),u(t))\,dt\\
 \text{s.t.}&\quad 
\dot q=f(q)+A(q)u,\ \ q(0)=q_0,\ q(T)=q_T,\ \ d_i(q(t))\ge 0\ \forall t,\ i=1,\dots,m .
\end{aligned}
\end{equation}
Assume:
\begin{enumerate}[(H1)]
\item $U\subset\mathbb R^d$ is compact and convex, $u(\cdot)\in L^\infty([0,T];U)$.
\item $f,A$ are globally Lipschitz on $\mathbb R^n$; $L(\cdot,\cdot)$ is continuous and convex in $u$,
and there exist $c_0,c_1>0$ such that $L(q,u)\ge -c_0+c_1\|u\|^2$.
\item Each $d_i\in C^2(\mathbb R^n)$ and the feasible set of \cref{eq:OCP_hard} is nonempty.
Moreover, (\emph{Slater}) there exists a strictly feasible pair $(\bar q,\bar u)$ satisfying $d_i(\bar q(t))\ge \delta$
for all $t\in[0,T]$ and all $i$, for some $\delta>0$.
\item (\emph{Existence of an optimal solution}) \cref{eq:OCP_hard} admits at least one optimal pair $(q^\star,u^\star)$.
\end{enumerate}
For $\mu>0$, define the \emph{log-barrier relaxed} (unconstrained) problems
\begin{equation}
\begin{aligned}
\label{eq:OCP_barrier}
\min_{u(\cdot)}&\ J_\mu(u)
:=\int_0^T \Big( L(q(t),u(t)) + \mu\sum_{i=1}^m \alpha_i\,b(d_i(q(t)))\Big)\,dt \\
\text{s.t.}&\quad
\dot q=f(q)+A(q)u,\ q(0)=q_0,\ q(T)=q_T,
\end{aligned}
\end{equation}
where $\alpha_i>0$ are fixed weights and $b(s)=-\log s$ (so \cref{eq:OCP_barrier} is defined only on trajectories with
$d_i(q(t))>0$ for all $t$). Let $(q^\mu,u^\mu)$ be a minimizer of \cref{eq:OCP_barrier}.

Then:
\begin{enumerate}[(a)]
\item (\emph{Strict feasibility}) For every $\mu>0$, $d_i(q^\mu(t))>0$ for all $t\in[0,T]$ and all $i$.
\item (\emph{Subsequence convergence}) For any sequence $\mu_k\downarrow 0$, there exists a subsequence (not relabeled) and
a feasible pair $(q^0,u^0)$ for \cref{eq:OCP_hard} such that
\[
q^{\mu_k}\to q^0 \ \text{uniformly on }[0,T],
\qquad
u^{\mu_k}\rightharpoonup^\ast u^0 \ \text{in }L^\infty([0,T];\mathbb R^d).
\]
\item (\emph{Optimality of cluster points}) Any cluster point $(q^0,u^0)$ obtained in (b) is optimal for \cref{eq:OCP_hard}, and
the optimal values satisfy $\lim_{\mu\downarrow 0}\inf J_\mu = \inf J$.
\end{enumerate}

Moreover, assume in addition:
\begin{enumerate}[(H5)]
\item (\emph{Uniform multiplier tightness}) The measures $\mu_i^\mu(dt):=\lambda_i^\mu(t)\,dt$ defined by
\begin{equation}
\label{eq:lambda_def}
\lambda_i^\mu(t):=\mu\,\alpha_i\,\frac{1}{d_i(q^\mu(t))}\ge 0
\quad\Longleftrightarrow\quad
\mu\,\alpha_i\,\nabla b(d_i(q^\mu(t))) = -\lambda_i^\mu(t)\,\nabla d_i(q^\mu(t))
\end{equation}
have uniformly bounded total mass:
$\sup_{\mu>0}\mu_i^\mu([0,T])<\infty$ for each $i$.
\end{enumerate}
Then, along the subsequence in (b), $\mu_i^{\mu_k}\overset{\ast}{\rightharpoonup} \mu_i$ as finite Radon measures,
and the limiting pair $(q^0,u^0)$ admits a costate $p$ and measures $\{\mu_i\}$ satisfying the state-constrained PMP in measure form.
In particular, complementarity holds:
\begin{equation}
\label{eq:measure_complementarity_limit}
\int_{[0,T]} d_i(q^0(t))\,\mu_i(dt)=0,\qquad \mathrm{supp}(\mu_i)\subset\{t:\ d_i(q^0(t))=0\}.
\end{equation}
\end{proposition}

\begin{proof}
\textbf{Step 1 (Existence and strict feasibility for the barrier problems).}
Fix $\mu>0$. By (H3) there exists at least one strictly feasible trajectory, hence $J_\mu$ is finite on a nonempty set.
Since $b(s)=-\log s$ diverges to $+\infty$ as $s\downarrow 0$ and is undefined for $s\le 0$,
any admissible minimizer must satisfy $d_i(q^\mu(t))>0$ for all $t$ and $i$, proving (a).
Existence of a minimizer for \cref{eq:OCP_barrier} follows from standard deterministic OCP existence arguments:
(H1) provides weak-$^\ast$ compactness of controls in $L^\infty$; (H2) gives existence/uniqueness and continuous dependence
of $q$ on $u$ (via global Lipschitz dynamics), and lower semicontinuity of the integral cost under weak-$^\ast$ convergence
(using convexity in $u$). Thus a minimizing sequence has a convergent subsequence whose limit attains the infimum.

\textbf{Step 2 (Compactness of $\{(q^\mu,u^\mu)\}$ as $\mu\downarrow 0$).}
Let $\bar u$ be the strictly feasible Slater control from (H3), and $\bar q$ its trajectory.
For any minimizer $(q^\mu,u^\mu)$,
\[
J_\mu(u^\mu)\le J_\mu(\bar u)=\int_0^T L(\bar q(t),\bar u(t))\,dt + \mu\sum_i \alpha_i\int_0^T b(d_i(\bar q(t)))\,dt
\le C_0 + \mu C_1,
\]
where $C_0,C_1<\infty$ are constants (since $d_i(\bar q(t))\ge\delta>0$ implies $b(d_i(\bar q(t)))$ is bounded).
By (H2), $L(q,u)\ge -c_0+c_1\|u\|^2$ yields a uniform $L^2$ bound on $u^\mu$ (hence also boundedness in $L^\infty$ because $u^\mu(t)\in U$).
Therefore $\{u^\mu\}$ is bounded in $L^\infty$, so by Banach--Alaoglu there exists a subsequence $\mu_k\downarrow 0$
with $u^{\mu_k}\rightharpoonup^\ast u^0$ in $L^\infty$.
By Lipschitz dynamics and bounded controls, the corresponding states $q^{\mu_k}$ are equi-Lipschitz and uniformly bounded,
so by Arzel\`a--Ascoli, $q^{\mu_k}\to q^0$ uniformly for some absolutely continuous $q^0$.
Passing to the limit in the ODE (continuous dependence of solutions on controls) gives that $q^0$ is the trajectory induced by $u^0$.
Since each $d_i(q^{\mu_k}(t))>0$ and $d_i$ is continuous, the uniform limit implies $d_i(q^0(t))\ge 0$ for all $t$,
so $(q^0,u^0)$ is feasible for \cref{eq:OCP_hard}. This proves (b).

\textbf{Step 3 (Value convergence and optimality of cluster points).}
We prove $\lim_{\mu\downarrow 0}\inf J_\mu = \inf J$ and that any cluster point is optimal.

\emph{(i) {Lim-inf} inequality.}
For the minimizing sequence $u^{\mu_k}$,
\[
\inf J_{\mu_k}=J_{\mu_k}(u^{\mu_k}) \ge \int_0^T L(q^{\mu_k}(t),u^{\mu_k}(t))\,dt = J(u^{\mu_k}),
\]
since the barrier term is nonnegative.
By convexity of $L$ in $u$ and uniform convergence $q^{\mu_k}\to q^0$, the integral functional $u\mapsto\int_0^T L(q[u](t),u(t))dt$
is weak-$^\ast$ lower semicontinuous, hence
\[
\liminf_{k\to\infty} \inf J_{\mu_k} \ \ge\ \liminf_{k\to\infty} J(u^{\mu_k})\ \ge\ J(u^0)\ \ge\ \inf J.
\]

\emph{(ii) {Lim-sup} inequality.}
Fix any feasible control $u$ for \cref{eq:OCP_hard} with trajectory $q$ (so $d_i(q(t))\ge 0$).
Let $(\bar q,\bar u)$ be the Slater strictly feasible pair with margin $\delta>0$.
Define a perturbed control
\[
u^{(\epsilon)} := (1-\epsilon)u+\epsilon \bar u,\qquad \epsilon\in(0,1).
\]
Then $u^{(\epsilon)}\in U$ a.e.\ by convexity of $U$.
Let $q^{(\epsilon)}$ be the induced trajectory.
By continuous dependence of the ODE solution on controls, $q^{(\epsilon)}\to q$ uniformly as $\epsilon\downarrow 0$.
Moreover, since $(\bar q,\bar u)$ is strictly feasible and $d_i$ is continuous, there exists $\epsilon_0>0$ such that
for all $\epsilon\in(0,\epsilon_0)$, $q^{(\epsilon)}$ is \emph{strictly feasible}, i.e.\ $\min_{t\in[0,T]} d_i(q^{(\epsilon)}(t))>0$
for all $i$ (intuitively: the perturbation moves the trajectory into the interior; formally, use uniform convergence plus
$d_i(\bar q(t))\ge\delta$ and continuity of $d_i$).
Hence $J_\mu(u^{(\epsilon)})$ is well-defined for all $\mu>0$.

Now choose a sequence $\epsilon=\epsilon(\mu)\downarrow 0$ such that
\[
\epsilon(\mu)\downarrow 0
\qquad\text{and}\qquad
\mu\,\log\!\big(1/\epsilon(\mu)\big)\to 0\quad(\mu\downarrow 0),
\]
e.g.\ $\epsilon(\mu)=\exp(-\sqrt{1/\mu})$.
Then $q^{(\epsilon(\mu))}\to q$ uniformly and $J(u^{(\epsilon(\mu))})\to J(u)$ by continuity of $L$ and dominated convergence theorem.
For the barrier term, strict feasibility implies $d_i(q^{(\epsilon(\mu))}(t))\ge c\,\epsilon(\mu)$ for some $c>0$
on $[0,T]$ (possibly after shrinking $\epsilon_0$); thus
\[
0\le \mu\alpha_i\int_0^T b(d_i(q^{(\epsilon(\mu))}(t)))\,dt
\ \le\ \mu\alpha_i T\big(\log(1/c)+\log(1/\epsilon(\mu))\big)\ \xrightarrow[\mu\downarrow 0]{}\ 0.
\]
Therefore $\limsup_{\mu\downarrow 0}\inf J_\mu \le \lim_{\mu\downarrow 0} J_\mu(u^{(\epsilon(\mu))}) = J(u)$.
Taking infimum over feasible $u$ yields $\limsup_{\mu\downarrow 0}\inf J_\mu \le \inf J$.

Combining (i) and (ii) gives $\lim_{\mu\downarrow 0}\inf J_\mu = \inf J$.
Finally, from the liminf chain $\inf J \le J(u^0)\le \liminf \inf J_{\mu_k}=\inf J$, we conclude $J(u^0)=\inf J$,
so $(q^0,u^0)$ is optimal for \cref{eq:OCP_hard}. This proves (c).

\textbf{Step 4 {(Adjoint equation in measure form and complementarity)}.}
For each $\mu$, the barrier problem \cref{eq:OCP_barrier} is unconstrained (aside from boundary conditions), hence admits the classical PMP.
The barrier contribution to the adjoint equation is
\[
-\nabla_q\Big(\mu\sum_i \alpha_ib(d_i(q^\mu(t)))\Big)
= -\sum_i \lambda_i^\mu(t)\,\nabla d_i(q^\mu(t)),
\qquad
\lambda_i^\mu(t)=\mu\alpha_i\frac{1}{d_i(q^\mu(t))}\ge 0,
\]
which is \cref{eq:lambda_def}.
Assumption (H5) implies the measures $\mu_i^\mu(dt):=\lambda_i^\mu(t)\,dt$ have uniformly bounded total mass, hence are weakly-$^\ast$
precompact in the space of finite Radon measures. Therefore, along a subsequence, $\mu_i^{\mu_k}\rightharpoonup^\ast \mu_i$.
Passing to the limit in the adjoint equation yields the measure-form adjoint relation of the state-constrained PMP, with multipliers $\{\mu_i\}$.

To prove complementarity \cref{eq:measure_complementarity_limit}, note that for each $\mu$,
\[
\int_0^T d_i(q^\mu(t))\,\mu_i^\mu(dt)
=\int_0^T d_i(q^\mu(t))\,\lambda_i^\mu(t)\,dt
=\int_0^T \mu\alpha_i\,dt
=\mu\alpha_i T \xrightarrow[\mu\downarrow 0]{} 0.
\]
Since $q^{\mu_k}\to q^0$ uniformly and $d_i$ is continuous, $d_i(q^{\mu_k})\to d_i(q^0)$ uniformly.
By weak-$^\ast$ convergence of measures, $\int d_i(q^{\mu_k})\,d\mu_i^{\mu_k}\to \int d_i(q^0)\,d\mu_i$.
Thus $\int d_i(q^0)\,d\mu_i=0$ and, since both terms are nonnegative, $\mu_i$ is supported on $\{t:\ d_i(q^0(t))=0\}$.
\end{proof}

We extend the \cref{prop:barrier_to_state_constrained_complete} in IPC barrier potential as well. The statement is addressed in \cref{cor:ipc_from_logbarrier}.

\begin{corollary}[IPC barrier inherits the log-barrier limit (primal) and contact-measure form (dual)]
\label{cor:ipc_from_logbarrier}
Consider the state-constrained OCP \cref{eq:OCP_hard} and assume the hypotheses of
Proposition~\ref{prop:barrier_to_state_constrained_complete} stated for the \emph{log-barrier} $b(s)=-\log s$.
Fix activation distances $\hat d_i>0$ and define the IPC barrier (clamped to zero outside the activation region) by
\begin{equation}
\label{eq:ipc_barrier_def}
b_{\mathrm{IPC}}(d;\hat d)
:=
\begin{cases}
(\hat d-d)^2\,\phi\!\big(\frac{d}{\hat d}\big), & 0<d<\hat d,\\
0, & d\ge \hat d,
\end{cases}
\qquad b(x)=-\log x .
\end{equation}
For $\mu>0$, consider the barrier-relaxed problems with running cost augmented by
\(
\mu\sum_i \alpha_i\, b_{\mathrm{IPC}}(d_i(q);\hat d_i)
\)
(where $\alpha_i>0$ are fixed weights and $d_i(q)>0$ is required whenever $b_{\mathrm{IPC}}$ is active).
Let $(q^\mu,u^\mu)$ be minimizers.

\smallskip
\noindent\textbf{(i) Primal limit.}
As $\mu\downarrow 0$, any sequence $\mu_k\downarrow 0$ admits a subsequence such that
\(
q^{\mu_k}\to q^\star
\)
(uniformly) and
\(
u^{\mu_k}\overset{\ast}{\rightharpoonup} u^\star
\),
where $(q^\star,u^\star)$ is an optimal solution of the original hard-constrained OCP \cref{eq:OCP_hard}.
Moreover, $\lim_{\mu\downarrow 0}\inf J_\mu=\inf J$.

\smallskip
\noindent\textbf{(ii) Multiplier/measure form (under the same tightness condition).}
Define the IPC multiplier densities on the active set $\{t:\ 0<d_i(q^\mu(t))<\hat d_i\}$ by
\begin{equation}
\label{eq:ipc_lambda_density}
\lambda_i^\mu(t)
:=\mu\,\alpha_i\Big[-\partial_d b_{\mathrm{IPC}}\big(d_i(q^\mu(t));\hat d_i\big)\Big]\ \ge 0,
\qquad
\mu_i^\mu(dt):=\lambda_i^\mu(t)\,dt .
\end{equation}
If $\sup_{\mu>0}\mu_i^\mu([0,T])<\infty$ (uniform tightness), then along a subsequence
$\mu_i^\mu \rightharpoonup^\ast \mu_i$ weakly-$^\ast$ as finite Radon measures, and the limiting necessary conditions
take the state-constrained PMP form with contact measures $\{\mu_i\}$.
In addition, the complementarity residual vanishes:
\begin{equation}
\label{eq:ipc_comp_residual}
\int_0^T d_i(q^\mu(t))\,\mu_i^\mu(dt)\ \xrightarrow[\mu\downarrow 0]{}\ 0,
\qquad\text{hence}\qquad
\mathrm{supp}(\mu_i)\subset\{t:\ d_i(q^\star(t))=0\}.
\end{equation}
\end{corollary}

\begin{proof}
\textbf{Step 1 (IPC barrier is sandwiched by log barriers on a near-contact region).}
Fix $i$ and write $x:=d/\hat d_i\in(0,1)$ on the active set. Since $0<d<\hat d_i$ implies
$0<(\hat d_i-d)^2\le \hat d_i^2$, we have the pointwise upper bound
\begin{equation}
\label{eq:ipc_upper_log}
0\le b_{\mathrm{IPC}}(d;\hat d_i)=(\hat d_i-d)^2\,b(x)\ \le\ \hat d_i^{\,2}\,b(x).
\end{equation}
On the \emph{near-contact} region $0<d\le \hat d_i/2$ we also have $(\hat d_i-d)^2\ge (\hat d_i/2)^2$, hence
\begin{equation}
\label{eq:ipc_lower_log}
b_{\mathrm{IPC}}(d;\hat d_i)\ \ge\ (\hat d_i/2)^{2}\,b(x),
\qquad 0<d\le \hat d_i/2.
\end{equation}
Thus, on any trajectory segment where contacts are potentially active, IPC is equivalent (up to constants) to a log barrier.

\textbf{Step 2 (Primal convergence via Proposition~\ref{prop:barrier_to_state_constrained_complete}).}
Consider the IPC-relaxed objective
\[
J_\mu^{\mathrm{IPC}}(u)=\int_0^T\Big(L(q,u)+\mu\sum_i \alpha_i\,b_{\mathrm{IPC}}(d_i(q);\hat d_i)\Big)\,dt.
\]
By \cref{eq:ipc_upper_log}, for any admissible $(q,u)$,
\[
J_\mu^{\mathrm{IPC}}(u)\ \le\ \int_0^T\Big(L(q,u)+\mu\sum_i \alpha_i\,\hat d_i^{\,2}\,\phi\!\Big(\frac{d_i(q)}{\hat d_i}\Big)\Big)\,dt
=:J_\mu^{\mathrm{log},U}(u),
\]
i.e., IPC is no more penalizing than a scaled log barrier.
Conversely, by \cref{eq:ipc_lower_log}, on times when $d_i(q)\le \hat d_i/2$ the IPC term dominates a (smaller) scaled log barrier.
This is sufficient for the standard barrier homotopy argument used in Proposition~\ref{prop:barrier_to_state_constrained_complete}:
(i) strict feasibility holds for each $\mu$ since $b_{\mathrm{IPC}}(d)\to\infty$ as $d\downarrow 0$;
(ii) the recovery sequence can be constructed exactly as in Proposition~\ref{prop:barrier_to_state_constrained_complete} by choosing
$\epsilon(\mu)\downarrow 0$ such that $\mu\,\hat d_i^{\,2}\,b(c\,\epsilon(\mu))\to 0$ (e.g.\ $\epsilon(\mu)=\exp(-1/\sqrt{\mu})$),
and then applying \cref{eq:ipc_upper_log} to bound the IPC penalty by the scaled log penalty.
Hence $\lim_{\mu\downarrow 0}\inf J_\mu^{\mathrm{IPC}}=\inf J$ and any cluster point of minimizers solves the hard-constrained OCP.
This proves (i).

\textbf{Step 3 {(Adjoint equation in measure form and complementarity)}.}
For each $\mu>0$ the IPC-relaxed problem is unconstrained (on the domain $d_i(q)>0$), so its PMP adjoint equation includes the smooth force
$-\mu\alpha_i\,\partial_d b_{\mathrm{IPC}}(d_i(q^\mu);\hat d_i)\,\nabla d_i(q^\mu)$.
Defining $\lambda_i^\mu$ by \cref{eq:ipc_lambda_density} yields the absolutely continuous measures $\mu_i^\mu(dt)=\lambda_i^\mu(t)\,dt$.
Under the same tightness condition as in Proposition~\ref{prop:barrier_to_state_constrained_complete}, weak-$^\ast$ sequential compactness of
$\mathcal M_+([0,T])$ implies $\mu_i^\mu\rightharpoonup^\ast \mu_i$ along a subsequence, and passing to the limit in the adjoint equation gives
the state-constrained measure form.

To show \cref{eq:ipc_comp_residual}, note that on the active set $0<d<\hat d_i$ we can write
$b_{\mathrm{IPC}}(d;\hat d_i)=\hat d_i^{\,2}(1-x)^2b(x)$ with $x=d/\hat d_i$.
A direct differentiation gives
\(
-d\,\partial_d b_{\mathrm{IPC}}(d;\hat d_i)
=
\hat d_i^{\,2}\,x(1-x)\big(2(1-x)b(x)+1\big),
\)
and using the elementary bound $xb(x)=x(-\log x)\le 1/e$ for $x\in(0,1]$ yields
\[
0\le -d\,\partial_d b_{\mathrm{IPC}}(d;\hat d_i)\ \le\ \hat d_i^{\,2}\Big(1+\frac{2}{e}\Big).
\]
Therefore,
\begin{equation}
\begin{aligned}
\int_0^T d_i(q^\mu(t))\,\mu_i^\mu(dt)
&= \mu\,\alpha_i\int_0^T \Big(-d_i(q^\mu(t))\,\partial_d b_{\mathrm{IPC}}(d_i(q^\mu(t));\hat d_i)\Big)\,dt \\
&\leq \mu\,\alpha_i\,T\,\hat d_i^{\,2}\Big(1+\frac{2}{e}\Big)\ \xrightarrow[\mu\downarrow 0]{} 0,        
\end{aligned}
\end{equation}
proving the vanishing complementarity residual and hence the support statement in \cref{eq:ipc_comp_residual}.
\end{proof}




\subsection{Inverse Hamiltonian design by potential reshaping}
\label{app:inverse_design}


We provide additional propositions and theorems to motivate our choice of Hamiltonian function and Hamiltonian dynamics. The optimally controlled dynamics admits a PMP Hamiltonian which solves a Two Point Boundary Value Problem (\textbf{TVBVP}). Therefore, it is worthwhile to state reversibility of the Hamiltonian flow and how it relates our choice of Hamiltonian function search space.

\begin{proposition}[Symplectic flow and the shooting map]
\label{lem:symplectic_flow_shooting}
Fix $\mathcal E$ and let $H(\cdot;\mathcal E)\in C^2(T^{*}\mathcal{Q})$.
Let $\Phi_H^t(\cdot;\mathcal E):T^{*}\mathcal{Q}\to T^{*}\mathcal{Q}$ be the Hamiltonian flow of $\dot x=X_H(x)$.
For any $(q_0,p_0)\in T^{*}{\mathcal{Q}}$ and any $t$ in the existence interval, the endpoint map
\[
(q(t),p(t))=\Phi_H^t(q_0,p_0;\mathcal E)
\]
is well-defined, and $\Phi_H^t(\cdot;\mathcal E)$ is a diffeomorphism with inverse $\Phi_H^{-t}(\cdot;\mathcal E)$.
In particular, the map
\[
\mathcal S_{\mathcal E}(p_0)\ :=\ \pi_Q\!\big(\Phi_H^{T}(q_0,p_0;\mathcal E)\big)\in Q
\]
(where $\pi_Q(q,p)=q$) is $C^1$ and encodes the \emph{shooting formulation} of the two-point boundary value problem.
\end{proposition}

\begin{proposition}[Local feed-forward representation of the TPBVP via an invertible shooting map]
\label{prop:tpbvp_feedforward_by_ift}
Consider the fixed-endpoint TPBVP
\begin{equation}
\label{eq:tpbvp_fixed_endpoint}
\dot q=\nabla_p H(q,p;\mathcal E),\qquad
\dot p=-\nabla_q H(q,p;\mathcal E),\qquad
q(0)=q_0,\ \ q(T)=q_T.
\end{equation}




Assume there exists a solution $(q^\star(\cdot),p^\star(\cdot))$ and define $p_0^\star:=p^\star(0)$.
Suppose the shooting map $\mathcal S_{\mathcal E}$ from Lemma~\ref{lem:symplectic_flow_shooting} satisfies the rank condition
\begin{equation}
\label{eq:rank_condition_shooting}
D\mathcal S_{\mathcal E}(p_0^\star)\ \in\ \mathbb R^{n\times n}\ \text{is nonsingular}\qquad (n=\dim Q).
\end{equation}
Then, by the implicit function theorem, there exist neighborhoods $\mathcal U$ of $q_T$ and $\mathcal V$ of $p_0^\star$
and a $C^1$ map
\[
\sigma_{\mathcal E}:\mathcal U\to \mathcal V
\quad\text{s.t.}\quad
\mathcal S_{\mathcal E}(\sigma_{\mathcal E}(q_T))=q_T,\ \ \sigma_{\mathcal E}(q_T)=p_0^\star.
\]
Consequently, the TPBVP \cref{eq:tpbvp_fixed_endpoint} is locally equivalent to the following \emph{feed-forward IVP}:
\[
(q(t),p(t))=\Phi_H^{t}\big(q_0,\sigma_{\mathcal E}(q_T);\mathcal E\big),\qquad t\in[0,T].
\]
That is, solving the TPBVP reduces (locally) to applying the feed-forward map $p_0=\sigma_{\mathcal E}(q_T)$ and integrating forward.
\end{proposition}

\paragraph{What does the rank condition mean}
Condition \eqref{eq:rank_condition_shooting} is precisely the \emph{nondegeneracy} that makes the endpoint map locally invertible.
It can be interpreted as local surjectivity of the map $p_0\mapsto q(T)$ around $p_0^\star$; equivalently,
small changes in the initial costate can steer the final position in all directions.
Without \eqref{eq:rank_condition_shooting}, the TPBVP may admit multiple solutions or no locally unique solution, and no smooth
feed-forward map $q_T\mapsto p_0$ can be asserted.

\begin{proposition}[Forward-only reduction via a costate feedback (Hamilton--Jacobi condition)]
\label{prop:tpbvp_forward_only_hj}
Assume there exists a $C^1$ function $W:[0,T]\times Q\to\mathbb R$ such that the \emph{feedback costate}
\[
p=\kappa(t,q):=\nabla_q W(t,q)
\]
satisfies, along the induced trajectory, the identity
\begin{equation}
\label{eq:hj_identity}
\partial_t W(t,q)\ +\ H\big(q,\nabla_q W(t,q);\mathcal E\big)\ =\ 0,
\qquad
W(T,q)=g(q)
\end{equation}
for some terminal cost $g$ (or more generally a terminal condition encoding the endpoint constraint).
Then the optimal phase trajectory satisfying the PMP conditions can be generated by the \emph{forward-only} closed-loop ODE
\[
\dot q(t)=\nabla_p H\big(q(t),\nabla_q W(t,q(t));\mathcal E\big),\qquad q(0)=q_0,
\]
with the costate recovered pointwise by $p(t)=\nabla_q W(t,q(t))$.
\end{proposition}

\paragraph{ Interpretation for Geometric reinforcement learning}
Proposition~\ref{prop:tpbvp_feedforward_by_ift} shows a minimal requirement for a feed-forward representation of a TPBVP:
a (local) inverse of the shooting map exists under a rank condition.
Proposition~\ref{prop:tpbvp_forward_only_hj} is stronger: it reduces the forward--backward system to a forward IVP once a costate
feedback $p=\kappa(t,q)$ exists, which is equivalent to the existence of a generating function $W$ solving a Hamilton--Jacobi equation.
In practice, a learned meta-map $\mathcal E\mapsto H(\cdot;\mathcal E)$ together with either (i) a learned initializer
$\sigma_{\mathcal E}$ or (ii) a learned costate feedback $\kappa_\xi(t,q;\mathcal E)$ provides a feed-forward surrogate of the TPBVP.

\begin{lemma}[Time-reversal symmetry for mechanical Hamiltonians]
\label{lem:mech_time_reversal}
Let $H(q,p)=K(p)+V(q)$ on $T^{*}\mathcal{Q}$, where $K\in C^2$ satisfies $K(-p)=K(p)$. Let $\Phi_H^t$ be the Hamiltonian flow and define the involution $\iota_R(q,p):=(q,-p)$.
Then, for all $t$ for which both sides are defined,
\[
\iota_R\circ \Phi_H^{t} \;=\; \Phi_H^{-t}\circ \iota_R.
\]
Equivalently, if $x(t)=(q(t),p(t))$ solves Hamilton's equations forward in time, then $\tilde x(t):=\iota_R(x(-t))=(q(-t),-p(-t))$ is also a forward-time solution.
\end{lemma}

\begin{proof}
Hamilton's equations are $\dot q=\nabla_p K(p)$ and $\dot p=-\nabla_q V(q)$.
Let $\tilde q(t):=q(-t)$ and $\tilde p(t):=-p(-t)$.
Then $\dot{\tilde q}(t)=-\dot q(-t)=-\nabla_p K(p(-t))=\nabla_p K(-p(-t))=\nabla_p K(\tilde p(t))$,
where we used $\nabla_p K(-p)=-\nabla_p K(p)$.
Also $\dot{\tilde p}(t)=-(-\dot p(-t))=\dot p(-t)=-\nabla_q V(q(-t))=-\nabla_q V(\tilde q(t))$. Thus $(\tilde q,\tilde p)$ satisfies Hamilton's equations, which is equivalent to the flow identity.
\end{proof}

\paragraph{ Feed-forward Hamiltonian and non-symplectic Hamiltonian dynamics}
Lemma~\ref{lem:mech_time_reversal} applies only to \emph{conservative} Hamiltonian dynamics (purely symplectic flows) which can reverse the p-dynamics given the terminal state of $p(T)$. Our deployed rollouts may include friction, barrier shaping, or other nonconservative corrections, in which case the induced
time-$\Delta t$ map is generally \emph{not} symplectic and the time-reversal identity may not hold.
In this work we use the conservative time-reversibility property only as \emph{motivation} for adopting a forward,
energy-shaped surrogate that is well-posed as an IVP and admits
Lyapunov-type descent guarantees. Our policy learner (\cref{eq:module_score_field}) learns a generalized Hamiltonian dynamics where dissipation and open-loop source term are considered as well. Formal approximability between such surrogates and the original PMP and TPBVP are considered as further research interests. Below we state our preliminary theorems.

\begin{theorem}[Local existence of a costate initializer map]
\label{thm:ift_costate}
Assume: (i) for a given $(q_0,q_T,\mathcal E)$ there exists a solution $(q^\star,p^\star)$ of \cref{eq:tpbvp_fixed_endpoint}
with $q^\star(0)=q_0$ and $q^\star(T)=q_T$;
(ii) the Jacobian $\nabla_{p_0}S(p_0^\star;q_0,q_T,\mathcal E)$ is nonsingular.
Then there exists a neighborhood $\mathcal U$ of $(q_0,q_T,\mathcal E)$ and a unique $C^1$ map
\[
\Pi^\star:\ (q_0,q_T,\mathcal E)\mapsto p_0^\star=\Pi^\star(q_0,q_T,\mathcal E)
\]
such that $S(\Pi^\star(q_0,q_T,\mathcal E);q_0,q_T,\mathcal E)=0$ for all $(q_0,q_T,\mathcal E)\in\mathcal U$.
\end{theorem}
\emph{Proof.} $S(\cdot;q_0,q_T,\mathcal E)$ is $C^1$ in $p_0$ under the $C^1$ regularity of $H$ and smooth dependence of ODE solutions on initial conditions.
Apply the implicit function theorem to $S(p_0;\,q_0,q_T,\mathcal E)=0$ at $(p_0^\star;q_0,q_T,\mathcal E)$. $\square$

\begin{corollary}[Forward-only reconstruction of the optimal trajectory]
\label{cor:forward_only_det}
Under the assumptions of Theorem~\ref{thm:ift_costate}, if a feed-forward model $\widehat \Pi_\xi(q_0,q_T,\mathcal E)$
approximates $\Pi^\star$ on $\mathcal U$, then the forward integration of \cref{eq:tpbvp_fixed_endpoint} initialized at
$(q_0,\widehat \Pi_\xi(q_0,q_T,\mathcal E))$ produces a trajectory whose terminal error satisfies
\[
\|q(T;q_0,\widehat \Pi_\xi,\mathcal E)-q_T\|
\le L_S\,\|\widehat \Pi_\xi(q_0,q_T,\mathcal E)-\Pi^\star(q_0,q_T,\mathcal E)\|
\]
for a local Lipschitz constant $L_S$, and the induced control $u^\star(q(t),p(t),t;\mathcal E)$ is near-optimal
whenever the resulting TPBVP residual is small (by standard stability of convex optimal control problems).
\end{corollary}

\subsection{Identifiability and convergence of Hamiltonian updates}
\label{app:learnability}

\paragraph{Discrete-time deterministic setting.}
Fix an environment $\mathcal E$ and a time step $\Delta t>0$. Let $(q_t,p_t)_{t=0}^{T}$ denote a (measured or high-fidelity)
trajectory generated by an (unknown) optimal control law. Assume the dynamics class and the control penalty are fixed/known, and
we learn only the \emph{potential weights} in a reduced Hamiltonian of the form
\begin{equation}
\label{eq:H_eta_discrete}
H_\eta(q,p,t;\mathcal E)
:=  \frac{1}{2}p^TM^{-1}p
+\mathcal R(q, t;\eta_{\xi}(\mathcal{E}))
\end{equation}
where $R(q, t;\eta_{\xi}(\mathcal{E})):=\eta^T\phi(q,t;\mathcal E)$ refer fixed-form potential functionals (determined from context of $\mathcal{E}$, and one example is \cref{eq:module_potentials}) and
$\eta:= \eta_{\xi}(\mathcal{E})\in\mathbb R^{m}$ are the weights (optionally constrained $\eta\ge 0$ via a parametrization).
We fix a gauge by assuming $\phi$ contains \emph{no pure constant feature} (or, equivalently, we impose a normalization constraint
such as $\int \mathcal R(q)\,dq=0$); this removes the additive-constant ambiguity.

\paragraph{Observation model: identifying $\eta$ from discrete costate increments.}
Assume the symplectic Euler form
\begin{equation}
\begin{aligned}
\label{eq:disc_costate_model}
p_{t+1}&=p_t-\Delta t~\nabla_q H_\eta(q_t,p_t,t;\mathcal E)+\varepsilon_t, \\
q_{t+1}&= q_t + \Delta t M^{-1} p_{t+1},\quad t=0,\dots,T-1, 
\end{aligned}
\end{equation}
where $\varepsilon_t$ collects discretization/modeling error ($\varepsilon_t=0$ in the noiseless case).
Since only $\mathcal R_\eta$ depends on $\eta$ and $\nabla_q \mathcal R_\eta=\sum_{j=1}^m \eta_j \nabla_q\phi_j$,
define the regression target
\begin{equation}
\label{eq:regression_target}
y_t
:=\frac{p_t-p_{t+1}}{\Delta t}\;=\;\sum_{j=1}^m \eta_j~g_{t,j}\;+\;\tilde\varepsilon_t,
\qquad
g_{t,j}:=\nabla_q\phi_j(q_t,t;\mathcal E),
\end{equation}
with $\tilde\varepsilon_t:=-\varepsilon_t/\Delta t$.
Stacking $y_t\in\mathbb R^{d_q}$ gives a linear least-squares problem in $\eta$.

\begin{assumption}[Persistent excitation on visited paths]
\label{ass:PE}
Let $G\in\mathbb R^{m\times m}$ be the (trajectory) Gram matrix
\[
G:=\sum_{t=0}^{T-1} \mathbf G_t,\qquad
\mathbf G_t:=\Big[\,\langle g_{t,i},g_{t,j}\rangle\,\Big]_{i,j=1}^m,
\]
where $\langle\cdot,\cdot\rangle$ is the Euclidean inner product in $\mathbb R^{d_q}$.
Assume $G\succeq \gamma I$ for some $\gamma>0$.
\end{assumption}

\begin{theorem}[Identifiability of potential weights from deterministic discrete paths]
\label{thm:identifiability_discrete}
Consider the noiseless case $\varepsilon_t\equiv 0$ in \cref{eq:disc_costate_model}.
Under Assumption~\ref{ass:PE}, the weights $\eta$ in \cref{eq:H_eta_discrete} are uniquely determined by the observed path
$\{(q_t,p_t)\}_{t=0}^{T}$ through \cref{eq:regression_target}; i.e., if $\eta$ and $\eta'$ yield the same
$\{y_t\}_{t=0}^{T-1}$, then $\eta=\eta'$.
\end{theorem}

\begin{proof}
If $\eta$ and $\eta'$ yield the same $y_t$, then $\sum_{j=1}^m(\eta_j-\eta'_j)g_{t,j}=0$ for all $t$.
Let $\delta:=\eta-\eta'$. Taking inner products with $g_{t,i}$ and summing over $t$ gives
$(G\delta)_i=\sum_{t}\sum_{j}\delta_j\langle g_{t,i},g_{t,j}\rangle=0$ for all $i$, hence $G\delta=0$.
Assumption~\ref{ass:PE} implies $\delta=0$, so $\eta=\eta'$.
\end{proof}

\paragraph{Learning objective and Hamiltonian update convergence in learning time.}
Given a dataset of trajectories (possibly multiple rollouts and/or multiple environments), define the empirical loss
\begin{equation}
\label{eq:ls_loss_eta}
\mathcal J(\eta)
:=\frac12\sum_{t=0}^{T-1}\big\|y_t-\sum_{j=1}^m \eta_j g_{t,j}\big\|_2^2
+\frac{\lambda}{2}\|\eta\|_2^2,
\qquad \lambda\ge 0,
\end{equation}
where $y_t$ and $g_{t,j}$ are defined by \cref{eq:regression_target}. This is a convex quadratic in $\eta$.

\begin{theorem}[Strong convexity, convergence of $\eta_k$, and vanishing Hamiltonian updates]
\label{thm:eta_convergence_dH}
Assume either (i) $\lambda>0$, or (ii) $\lambda=0$ and Assumption~\ref{ass:PE} holds.
Then $\mathcal J$ is $\mu$-strongly convex for some $\mu>0$ and has a unique minimizer $\eta^\star$.
Consider gradient descent in learning time $k$:
\[
\eta_{k+1}=\eta_k-\alpha\nabla \mathcal J(\eta_k),
\qquad 0<\alpha<\frac{2}{L},
\]
where $L$ is the Lipschitz constant of $\nabla\mathcal J$.
Then $\eta_k\to\eta^\star$ linearly and the induced Hamiltonian updates vanish on any bounded set:
for any compact $\mathcal K\subset Q\times \mathbb R^{d_q}$,
\begin{equation}
\label{eq:dH_vanish}
\sup_{(q,p)\in\mathcal K}\big|H_{\eta_{k+1}}(q,p,t;\mathcal E)-H_{\eta_k}(q,p,t;\mathcal E)\big|
\ \xrightarrow[k\to\infty]{}\ 0.
\end{equation}
Moreover, along any observed trajectory points $\{q_t\}$,
\begin{equation}
\label{eq:dH_rate}
\max_{t\in\{0,\dots,T\}}\big|H_{\eta_{k+1}}(q_t,p_t,t;\mathcal E)-H_{\eta_k}(q_t,p_t,t;\mathcal E)\big|
\le \|\eta_{k+1}-\eta_k\|_2 \cdot \max_t \|\phi(q_t,t;\mathcal E)\|_2
\ \xrightarrow[k\to\infty]{}\ 0.
\end{equation}
\end{theorem}

\begin{proof}
Under the stated conditions, the Hessian of $\mathcal J$ is $G+\lambda I\succeq \mu I$ for some $\mu>0$,
so $\mathcal J$ is strongly convex with unique minimizer $\eta^\star$. Standard gradient-descent theory (c.f. \citet{NocedalWright2006NumericalOptimization}) yields linear convergence
$\|\eta_k-\eta^\star\|\le \rho^k\|\eta_0-\eta^\star\|$ for some $\rho\in(0,1)$.
Since $H_\eta$ depends on $\eta$ only through $\mathcal R(q,t;\mathcal E)=\eta^\top\phi(q,t;\mathcal E)$,
we have the Lipschitz relation
$|H_{\eta'}-H_\eta|=|(\eta'-\eta)^\top\phi|\le \|\eta'-\eta\|\cdot \|\phi\|$.
Combining with $\|\eta_{k+1}-\eta_k\|\to 0$ yields \cref{eq:dH_vanish}--\cref{eq:dH_rate}.
\end{proof}

\begin{remark}[Interpreting $dH/dk$ (learning-time derivative)]
A convenient discrete proxy is
\[
\frac{d}{dk}H_{\eta_k}(q,p)\;\approx\;H_{\eta_{k+1}}(q,p)-H_{\eta_k}(q,p)
=\phi(q,t;\mathcal E)^\top(\eta_{k+1}-\eta_k)
=-\alpha~\phi(q,t;\mathcal E)^\top\nabla\mathcal J(\eta_k).
\]
Thus, once the learning dynamics drive $\|\nabla\mathcal J(\eta_k)\|\to 0$, the Hamiltonian update magnitude
vanishes uniformly on compact sets where $\phi$ is bounded.
\end{remark}

\begin{remark}[Multiple environments and environment-conditioned weights]
If $\eta$ is environment-dependent (e.g.\ $\eta=\eta(\mathcal E)$), the identifiability and strong convexity conditions
apply \emph{per environment} with the corresponding Gram matrix $G(\mathcal E)$ built from paths sampled in that environment.
If $\eta_\xi(\mathcal E)$ is produced by a neural network, the optimization is nonconvex in $\xi$; however, the same Gram
condition provides a \emph{local} identifiability statement for the induced weights and implies vanishing Hamiltonian updates
whenever the learning iterates converge to a stationary point.
\end{remark}
\subsection{Coupled stagewise goal selection and energy {reshaping}}
\label{app:coupled_stagewise}


\paragraph{Coupled stagewise problem: goal selection and Hamiltonian {reshaping}.}
In a stagewise motion planner, each local OCP (stage $\tau$) requires a terminal goal
$q_{\mathrm{target}}(\mathcal{E}_{\tau})$, while the next stage’s initial condition is inherited from the current stage’s terminal state.
Hence, if the terminal goal is produced by a policy, the overall procedure becomes \emph{coupled}:
the goal policy determines $q_{\mathrm{target}}(\mathcal{E}_{\tau})$, but its consequence is only revealed after integrating the stage OCP.

To formalize this coupling, let $x_\tau=(q_\tau,p_\tau)$ denote the augmented state, and assume the stagewise dynamics are induced
by a barrier-shaped reduced Hamiltonian family $H_{\eta_\tau}$ with weights $\eta_\tau=\eta_\xi(\mathcal E_\tau)$ and fixed horizon $T_\tau$.
Let $\Phi^{T_\tau}_{\eta_\tau}$ be the flow map at stage $\tau$. Let a goal policy $\pi_\gamma$ output the next stage terminal goal as
\begin{equation}
q_{\mathrm{target}}(\mathcal{E}_\tau) = \pi_\gamma(\mathcal E_\tau, x_\tau, \mathcal M_\tau),
\label{eq:goal_policy}
\end{equation}
where $\mathcal M_\tau$ is any memory/context (e.g.\ visited-map encoding, stage index, or recurrent hidden state).
Given $(x_\tau,q_{\mathrm{target}}(\mathcal{E}_{\tau}))$, the inverse-design reshaping step selects a weight vector $\eta_\tau$ (or equivalently $\xi$) so that the Hamiltonian rollout hits $q_{\mathrm{target}}(\mathcal{E}_{\tau})$.
Denote this selection operator by $\mathsf S$:
\begin{equation}
\eta_\tau = \mathsf S\big(x_\tau,q_{\mathrm{target}}(\mathcal{E}_{\tau}),\mathcal E_\tau\big),
\qquad\text{and}\qquad
x_{\tau+1} = \Phi^{T_\tau}_{\eta_\tau}\big(x_\tau;\mathcal E_\tau\big).
\label{eq:selection_and_transition}
\end{equation}

\paragraph{Composition yields a fixed-point (dueling) system.}
Substituting \cref{eq:goal_policy} into \cref{eq:selection_and_transition} yields the closed-loop stage map
\begin{equation}
x_{\tau+1}
=
\underbrace{\Phi^{T_\tau}_{\mathsf S(\,x_\tau,\pi_\gamma(\mathcal E_\tau,x_\tau,\mathcal M_\tau),\mathcal E_\tau\,)}
\big(x_\tau;\mathcal E_\tau\big)}_{=:~\mathcal T_{\gamma}(x_\tau,\mathcal M_\tau,\mathcal E_\tau)}.
\label{eq:closed_loop_stage_map}
\end{equation}
Thus the goal policy and the Hamiltonian-shaping inverse design are not independent modules; they form a coupled system
whose behavior is determined by the composite operator $\mathcal T_{\gamma}$.
The “dueling” phenomenon is precisely that $q_{\mathrm{target}}(\mathcal{E}_{\tau})$ is chosen before the OCP is solved, but its quality
can only be evaluated after the stage transition is realized. We state the following proposition as how to view sensor policy's role in the training pipeline but comment out that this joint learning problem remains an open question.

\begin{proposition}[Coupled stagewise inverse design as a fixed-point / bilevel problem]
\label{prop:coupled_goal_and_energy}
Assume: (i) the reshaping selection $\mathsf S$ is well-defined locally (e.g.\ by \cref{prop:tpbvp_feedforward_by_ift} with a
rank condition ensuring local solvability of $q(T)=q_{\mathrm{target}}(\mathcal{E}_{\tau})$ under fixed inherited $x_\tau$),
and (ii) the flow $\Phi^{T_\tau}_{\eta}$ depends smoothly on $\eta$ and $x_\tau$ on the region of interest.
Then learning a goal policy $\pi_\gamma$ jointly with energy weights (or their generator) is equivalently posed as either:

\smallskip
\noindent\textbf{(a) A fixed-point learning problem:} find $\gamma$ (and the reshaping module parameters) such that the induced
closed-loop stage map \cref{eq:closed_loop_stage_map} has desirable invariants (e.g.\ reaches a terminal set, avoids cycles,
satisfies progress constraints), i.e.\ $x_{\tau+1}=\mathcal T_\gamma(x_\tau,\mathcal M_\tau,\mathcal E_\tau)$ achieves a
prescribed global objective.

\smallskip
\noindent\textbf{(b) A bilevel optimization problem:} the goal policy is an outer decision variable that selects $q_{\mathrm{target}}(\mathcal{E}_{\tau})$,
while the inner level solves a stagewise inverse-design problem (reshaping) to realize that goal:
\begin{equation}  
\begin{aligned}
\min_{\gamma}&\ \mathbb E\Big[\sum_{\tau=0}^{S-1} \mathcal J_{\rm global}(x_\tau,q_{\mathrm{target}}(\mathcal{E}_{\tau}),\mathcal E_\tau)\Big]
\\
\text{s.t.}\quad&
q_{\mathrm{target}}(\mathcal{E}_{\tau})=\pi_\gamma(\cdot), \\
&\eta_\tau\in\arg\min_{\eta}\ \|\pi_q\Phi^{T_\tau}_{\eta}(x_\tau)-q_{\mathrm{target}}(\mathcal{E}_{\tau})\|^2+\mathrm{Reg}(\eta),
\\ 
&x_{\tau+1}=\Phi^{T_\tau}_{\eta_\tau}(x_\tau).    
\end{aligned}
\end{equation}
In either view, evaluating the quality of $q_{\mathrm{target}}(\mathcal{E}_{\tau})$ necessarily requires integrating (or rolling out) the
stage dynamics, so sequential solution is intrinsic to the problem structure.
\end{proposition}

\paragraph{Practical resolution: amortize the inner solve and train end-to-end}
\label{rem:amortize_inner}
In implementation, the inner ``reshaping'' solve (selecting $\eta_\tau$ to hit $q_{\mathrm{target}}(\mathcal{E}_{\tau})$) is typically amortized
by a feed-forward policy $\eta_\xi(\mathcal E_\tau,\cdot)$ and trained through differentiable rollouts.
This converts the bilevel structure into a single differentiable computation graph:
\[
x_{\tau+1} = \Phi^{T_\tau}_{\eta_\xi(\mathcal E_\tau,x_\tau,\mathcal M_\tau)}(x_\tau),
\qquad
q_{\mathrm{target}}(\mathcal{E}_{\tau})=\pi_\gamma(\mathcal E_\tau,x_\tau,\mathcal M_\tau),
\]
with gradients backpropagated through the unrolled stage map.


\paragraph{When sequential coupling is unavoidable}
\label{rem:sequential_unavoidable}
Even if $\pi_\gamma$ is deterministic, the mapping from $(\mathcal E_\tau,x_\tau)$ to the realized next state $x_{\tau+1}$
depends on the stagewise OCP solution operator. Therefore the next-stage goal cannot be certified without either
(i) executing the inner rollout/solve, or (ii) learning a sufficiently accurate surrogate for that solve.
This is a structural property of stagewise planning and not an artifact of a particular algorithm.


\section{Algorithm Implementation Details}
\label{app:algorithm_details}

This appendix gives the discrete-time implementation of the port--Hamiltonian navigation policy in
~\cref{eq:reduced_ham_flow_short} together with the module interface in
~\cref{eq:module_template_hamiltonian} and the dissipative/port form in
~\cref{eq:module_score_field}. We provide ~\cref{algo:online_grl_snam} (online GRL--SNAM)
with explicit equation references (no section cross-references).

Throughout, we use: environment state $\mathcal{E}_n$, active constraints $\mathcal{C}_n$, module
Hamiltonians $H_k$ and responses $\mathsf{R}_k$, navigator/meta-policy $g_\xi$, surrogate potential
$\mathcal{R}(q;\eta,\mathcal{E})$, and port--Hamiltonian correction parameters
$(\mu, u_f)$. The phase-space state is $z_n=(q_n,p_n)$.

\paragraph{Discrete-time convention.}
We use a fixed integration step $\tau>0$ and discrete index $n\in\mathbb{N}$ with physical time
$t_n := n\tau$. All horizons $(T_y,T_f,T_o)$ are integers measured in integration steps.

\paragraph{Configuration decomposition.}
When we need the robot center/position $\vec c_n$ from the configuration $q_n$, we write
\begin{equation}
\label{eq:impl:config_decomp}
q_n = (\vec c_n, \text{other coordinates}),\qquad \vec c_n := \mathrm{pos}(q_n).
\end{equation}

\paragraph{Active-set definition.}
Given signed distance functions $\{d_i(\cdot;\mathcal{E}_n)\}$ and activation threshold $\hat d$, the
active set used throughout is
\begin{equation}
\label{eq:impl:active_set}
\mathcal{C}_n := \{\, i \mid d_i(q_n;\mathcal{E}_n)\le \hat d \,\}.
\end{equation}

\paragraph{Surrogate energy and Hamiltonian.}
For $\eta=(\beta,\lambda,\{\alpha_i\}_{i\in\mathcal{C}_n})$, we assemble
\begin{align}
\label{eq:impl:surrogate_energy}
\mathcal{R}(q_n;\eta,\mathcal{E}_n)
&=
E_{\mathrm{sensor}}(q_n;\mathcal{E}_n,\omega_y)
+\beta\,E_{\mathrm{goal}}(q_n;\mathcal{E}_n,\omega_g) \\
&+\lambda\,E_{\mathrm{obj}}(q_n;\omega_d)
+\sum_{i\in\mathcal{C}_n}\alpha_i\, b(d_i(q_n;\mathcal{E}_n);\hat d,\omega_b),\\
\label{eq:impl:surrogate_hamiltonian}
H(q_n,p_n;\eta,\mathcal{E}_n)
&=\tfrac12 p_n^\top M^{-1}p_n+\mathcal{R}(q_n;\eta,\mathcal{E}_n).
\end{align}

\paragraph{Block-selectors for dissipation/ports (frame-only).}
To apply damping/ports only on the frame momentum coordinates, we implement
\begin{equation}
\label{eq:impl:block_selectors}
\Gamma(\mu) = \mathrm{blkdiag}(0,\Gamma_f(\mu),0),\qquad
G = \mathrm{blkdiag}(0,G_f,0),\qquad
\Gamma_f(\mu)=\mu I,\ \ G_f=I,
\end{equation}
and use $\Gamma(\mu),G$ in the discrete update below.

\paragraph{Gradients and one-step update (symplectic Euler).}
We compute
\begin{equation}
\label{eq:impl:gradients}
v_n := \nabla_p H(q_n,p_n)=M^{-1}p_n,\qquad
g_n := \nabla_q H(q_n,p_n)=\nabla_q \mathcal{R}(q_n;\eta,\mathcal{E}_n),
\end{equation}
and take the symplectic Euler discretization of ~\cref{eq:reduced_ham_flow_short}:
\begin{align}
\label{eq:impl:ph_euler_p}
p_{n+1} &= p_n-\tau\,g_n-\tau\,\Gamma(\mu)\,v_n+\tau\,G\,u_f,\\
\label{eq:impl:ph_euler_q}
q_{n+1} &= q_n+\tau\,M(q_n)^{-1}p_{n+1}.
\end{align}

\paragraph{Observables and target.}
We use
\begin{align}
\label{eq:impl:observables}
\mathrm{clr}_n &= \min\Big\{\min_i d_i(q_{n+1};\mathcal{E}_n),\ \min_{\ell,i} d_i(q_\ell^{(o)};\mathcal{E}_n)\Big\},\\
\mathrm{dist}_n &= \|\vec c_{n+1}-\mathbf{x}_g\|,\qquad
\mathrm{speed}_n = \|M(q_{n+1})^{-1}p_{n+1}\|,
\end{align}
and the observable vector
\begin{equation}
\label{eq:impl:observable_vector}
y_n=\begin{bmatrix}-\mathrm{clr}_n\\ \mathrm{dist}_n\\ -\mathrm{speed}_n\end{bmatrix}.
\end{equation}
We denote the (task-level) target as
\begin{equation}
\label{eq:impl:observable_target}
y_n^\star=\begin{bmatrix}-\mathrm{clr}^\star\\ \mathrm{dist}^\star\\ -\mathrm{speed}^\star\end{bmatrix},
\end{equation}
where $(\mathrm{clr}^\star,\mathrm{dist}^\star,\mathrm{speed}^\star)$ are fixed design setpoints.

\paragraph{Active parameter vector and Jacobian updates.}
Let $\mathcal{I}_n\subseteq \mathcal{C}_n$ be the subset of nearby obstacle indices used for the local
weight update, and define
\begin{equation}
\label{eq:impl:active_param_vector}
\tilde\eta_n=[\beta_n,\lambda_n,\{\alpha_{n,i}\}_{i\in\mathcal{I}_n}]^\top,\qquad
\zeta_n=[\tilde\eta_n,\mu_n]^\top.
\end{equation}
We maintain EMA-smoothed secant Jacobians
\begin{align}
\label{eq:impl:secant_J}
\widetilde{J}_n&=\frac{(y_n-y_{n-1})(\zeta_n-\zeta_{n-1})^\top}{\|\zeta_n-\zeta_{n-1}\|_2^2+\varepsilon},\qquad
J_n=\rho J_{n-1}+(1-\rho)\widetilde{J}_n,\\
\label{eq:impl:secant_P}
\widetilde{P}_n&=\frac{(y_n-y_{n-1})(u_{f,n}-u_{f,n-1})^\top}{\|u_{f,n}-u_{f,n-1}\|_2^2+\varepsilon},\qquad
P_n=\rho P_{n-1}+(1-\rho)\widetilde{P}_n.
\end{align}

\paragraph{Tikhonov LS update, projection, residual, and port correction.}
With desired increment $\Delta y_n^{\mathrm{des}}:=y_n^\star-y_n$, we solve
\begin{align}
\label{eq:impl:tikhonov_update}
\Delta\zeta_n&=(J_n^\top J_n+\lambda_\zeta I)^{-1}J_n^\top \Delta y_n^{\mathrm{des}},\qquad
\zeta_n^{\mathrm{upd}}=\Pi_{\mathbb{R}_+}\big(\zeta_n+\vec{\kappa}\odot\Delta\zeta_n\big),\\
\label{eq:impl:residual}
r_n&=\Delta y_n^{\mathrm{des}}-J_n\Delta\zeta_n,\\
\label{eq:impl:port_update}
u_{f,n}^{\mathrm{upd}}&=(P_n^\top P_n+\lambda_u I)^{-1}P_n^\top r_n,\qquad
u_{f,n}^{\mathrm{upd}}\leftarrow \mathrm{clip}(u_{f,n}^{\mathrm{upd}},\mathcal{U}).
\end{align}

\begin{algorithm}[htbp]
\caption{Online GRL--SNAM: Navigator-driven Hamiltonian composition with QoI feedback}
\label{algo:online_grl_snam}
\begin{algorithmic}[1]
\State \textbf{Input:} goal $\mathbf{x}_g$, initial $z_0=(q_0,p_0)$ (typically $p_0{=}0$), step $\tau$, horizons $(T_y,T_f,T_o)$, max steps $N_{\max}$
\State \textbf{Init:} $n\!\leftarrow\!0$, environment $\mathcal{E}_0$, active set $\mathcal{C}_0\!\leftarrow\!\emptyset$
\State \textbf{Init:} meta state $(\eta_0,\mu_0,u_{f,0})$ from $g_\xi$, Jacobians $J_0,P_0$ (~\cref{eq:impl:secant_J,eq:impl:secant_P})
\While{$\neg$\textsc{ReachedGoal}$(\vec{c}_n,\mathbf{x}_g)$ \textbf{and} $n<N_{\max}$} \Comment{$\vec c_n=\mathrm{pos}(q_n)$ from ~\cref{eq:impl:config_decomp}}
    \State \textbf{(A) Sensor refresh and active-set discovery (low rate):}
    \If{$n \equiv 0 \pmod{T_y}$}
        \State Query $\pi_y$ with $\mathcal{Q}_y^n=(\mathcal{E}_n,\mathcal{C}_n,z_{y,0}{=}z_n,H_y(\cdot;\eta_n),T_y;\theta_y)$
        \State Receive $\mathsf{R}_y^n=\{z_{0:T_y}^{(y)},s_{0:T_y}^{(y)},\Delta\mathcal{E}_n\}$
        \State Update $\mathcal{E}_n \leftarrow \mathcal{E}_n \cup \Delta\mathcal{E}_n$
        \State Update $\mathcal{C}_n \leftarrow \{i \mid d_i(q_n;\mathcal{E}_n)\le \hat d\}$ \Comment{~\cref{eq:impl:active_set}}
    \EndIf

    \State \textbf{(B) Frame and shape rollouts (medium/high rate):}
    \If{$n \equiv 0 \pmod{T_f}$}
        \State Query $\pi_f$ with $\mathcal{Q}_f^n=(\mathcal{E}_n,\mathcal{C}_n,z_{f,0}{=}z_n,H_f(\cdot;\eta_n,\mu_n,u_{f,n}),T_f;\theta_f)$
        \State Receive $\mathsf{R}_f^n=\{z_{0:T_f}^{(f)},s_{0:T_f}^{(f)},\{v_\ell^{(f)}\}_{\ell=0}^{T_f}\}$
    \EndIf
    \State Query $\pi_o$ with $\mathcal{Q}_o^n=(\mathcal{E}_n,\mathcal{C}_n,z_{o,0}{=}z_n,H_o(\cdot;\eta_n),T_o;\theta_o)$
    \State Receive $\mathsf{R}_o^n=\{z_{0:T_o}^{(o)},s_{0:T_o}^{(o)},\{\min_i d_i(q_\ell^{(o)};\mathcal{E}_n)\}_{\ell=0}^{T_o}\}$

    \State \textbf{(C) Meta-policy proposal (environment \& response $\to$ weights):}
    \State Build tokens $\mathcal{T}_n$ from $(\mathcal{E}_n,\mathcal{C}_n,z_n,\mathbf{x}_g,\mathsf{R}_f^n,\mathsf{R}_o^n)$
    \State $[\eta_n^{\mathrm{prop}},\mu_n^{\mathrm{prop}},u_{f,n}^{\mathrm{prop}}] \leftarrow g_\xi(\mathcal{T}_n;\xi)$
    \State \hspace{1em}with $\eta_n^{\mathrm{prop}}=(\beta_n,\lambda_n,\{\alpha_{n,i}\}_{i\in\mathcal{C}_n}) \in \mathbb{R}_+^{2+|\mathcal{C}_n|}$

    \State \textbf{(D) Compose surrogate Hamiltonian:} 
    % \Comment{~\cref{eq:impl:surrogate_energy,eq:impl:surrogate_hamiltonian}}
    \Statex \hspace{1em} Assemble $\mathcal{R}(q_n;\eta_n^{\mathrm{prop}},\mathcal{E}_n)$ via ~\cref{eq:impl:surrogate_energy}
    \Statex \hspace{1em} Assemble $H(q_n,p_n;\eta_n^{\mathrm{prop}},\mathcal{E}_n)$ via ~\cref{eq:impl:surrogate_hamiltonian}

    \State \textbf{(E) One-step port--Hamiltonian integration:}
    \State Compute $(v_n,g_n)$ via ~\cref{eq:impl:gradients}
    \State Apply frame-only dissipation/port via ~\cref{eq:impl:block_selectors} and step ~\cref{eq:impl:ph_euler_p,eq:impl:ph_euler_q}
    \State Set $z_{n+1}=(q_{n+1},p_{n+1})$

    \State \textbf{(F) Observable extraction:} \Comment{~\cref{eq:impl:observables,eq:impl:observable_vector}}
    \State Compute $(\mathrm{clr}_n,\mathrm{dist}_n,\mathrm{speed}_n)$ via ~\cref{eq:impl:observables}
    \State Set $y_n$ via ~\cref{eq:impl:observable_vector} and target $y_n^\star$ via ~\cref{eq:impl:observable_target}

    \State \textbf{(G) Online adaptation of active weights and friction:}
    \State Select nearby obstacle indices $\mathcal{I}_n \subseteq \mathcal{C}_n$
    \State Form $\zeta_n$ via ~\cref{eq:impl:active_param_vector}; update $J_n$ via ~\cref{eq:impl:secant_J}
    \State Solve $\Delta\zeta_n$ and project to $\zeta_n^{\mathrm{upd}}$ via ~\cref{eq:impl:tikhonov_update}

    \State \textbf{(H) Port correction from residual:}
    \State Compute residual $r_n$ via ~\cref{eq:impl:residual}
    \State Update $P_n$ via ~\cref{eq:impl:secant_P} and solve $u_{f,n}^{\mathrm{upd}}$ via ~\cref{eq:impl:port_update}

    \State \textbf{(I) Commit parameters for next step:}
    \State Unpack $\zeta_n^{\mathrm{upd}}\mapsto(\beta_n^{\mathrm{upd}},\lambda_n^{\mathrm{upd}},\{\alpha_{n,i}^{\mathrm{upd}}\}_{i\in\mathcal{I}_n},\mu_n^{\mathrm{upd}})$
    \State Set $\eta_{n+1}\leftarrow(\beta_n^{\mathrm{upd}},\lambda_n^{\mathrm{upd}},\{\alpha_{n+1,i}\}_{i\in\mathcal{C}_n})$ with $\alpha_{n+1,i}{=}\alpha_{n,i}^{\mathrm{upd}}$ for $i\in\mathcal{I}_n$ and unchanged otherwise
    \State Set $\mu_{n+1}\leftarrow \mu_n^{\mathrm{upd}}$, \quad $u_{f,n+1}\leftarrow u_{f,n}^{\mathrm{upd}}$
    \State $n\leftarrow n+1$
\EndWhile
\State \textbf{Return:} trajectory $\{z_n\}_{n=0}^{N}$ and parameter history $\{(\eta_n,\mu_n,u_{f,n})\}_{n=0}^{N}$
\end{algorithmic}
\end{algorithm}

\subsection{Initialization}
\label{subsec:initialization}

\paragraph{State initialization.}
Initialize phase-space coordinates $z_0=(q_0,p_0)$ and set $p_0=0$ unless otherwise stated. When needed,
extract $\vec c_0=\mathrm{pos}(q_0)$ using ~\cref{eq:impl:config_decomp}.

\paragraph{Step size and horizons.}
Choose $\tau \in [10^{-2},5{\times}10^{-2}]$ and integer horizons $(T_y,T_f,T_o)$ in steps.
Typical settings are $T_y{=}10$, $T_f{=}5$, $T_o{=}1$ (sensor refresh slowest, shape fastest),
matching the temporal hierarchy in Fig.~\ref{fig:temporal_scales}.

\paragraph{Maximum horizon.}
Let $N_{\max}$ be the maximum number of discrete steps (e.g., $5{,}000$--$10{,}000$), giving
$t_{\max}=N_{\max}\tau$.

\subsection{Parameter Structure and Initialization}
\label{subsec:parameter_structure}

\paragraph{Meta-policy output and parameterization.}
The navigator outputs
\[
g_\xi(\mathcal{T}_n;\xi)=\big[\eta_n^{\mathrm{prop}},\mu_n^{\mathrm{prop}},u_{f,n}^{\mathrm{prop}}\big],
\qquad
\eta_n^{\mathrm{prop}}=(\beta_n,\lambda_n,\{\alpha_{n,i}\}_{i\in\mathcal{C}_n})\in\mathbb{R}_+^{2+|\mathcal{C}_n|}.
\]
The intra-term parameters $\omega=(\omega_y,\omega_M,\omega_g,\omega_d,\omega_b)$ are fixed across environments;
only $(\eta,\mu,u_f)$ vary online.

\paragraph{Frame-only dissipation/ports.}
We enforce frame-only damping/ports using the block selectors in ~\cref{eq:impl:block_selectors},
and the discrete update in ~\cref{eq:impl:ph_euler_p,eq:impl:ph_euler_q}.

\paragraph{Jacobian initialization.}
Initialize $J_0$ and $P_0$ to small-norm matrices (e.g., $10^{-2}$ times identity where compatible), or
zeros. They are refined online via the EMA secant updates in ~\cref{eq:impl:secant_J,eq:impl:secant_P}.

\subsection{Termination Conditions}
\label{subsec:termination}

We terminate if any of the following holds:

\paragraph{Success (goal reached).}
$\|\vec{c}_n-\mathbf{x}_g\|<\epsilon_{\text{goal}}$, with $\vec c_n=\mathrm{pos}(q_n)$ from ~\cref{eq:impl:config_decomp}.

\paragraph{Timeout.}
$n\ge N_{\max}$.

\paragraph{Collision (optional hard-stop).}
$\min_i d_i(q_n;\mathcal{E}_n)<0$ (penetration).

\paragraph{Stuck (optional).}
$\|\vec{c}_n-\vec{c}_{n-W}\|<\epsilon_{\text{stuck}}$ for a fixed window $W$ (e.g., $W{=}50$ steps).

\subsection{Hierarchical Query Flow}
\label{subsec:hierarchical_queries}

GRL--SNAM uses the query--response interface consistent with ~\cref{eq:module_template_hamiltonian}:

\paragraph{Module query interface.}
Each module is queried with $\mathcal{Q}_k^n=(\mathcal{E}_n,\mathcal{C}_n,z_{k,0}{=}z_n,H_k(\cdot),T_k;\theta_k)$ and returns
\[
\mathsf{R}_k^n=\{z_{0:T_k}^{(k)},s_{0:T_k}^{(k)},\mathrm{QoI}_k\}.
\]
Concrete QoIs used in Algorithm~\cref{algo:online_grl_snam} are:
$\mathrm{QoI}_y=\Delta\mathcal{E}_n$,
$\mathrm{QoI}_f=\{v_\ell^{(f)}\}_{\ell=0}^{T_f}$, and
$\mathrm{QoI}_o=\{\min_i d_i(q_\ell^{(o)};\mathcal{E}_n)\}_{\ell=0}^{T_o}$.

\paragraph{Rate hierarchy.}
Sensor is queried every $T_y$ steps (active-set update via ~\cref{eq:impl:active_set}),
frame every $T_f$ steps (stage guidance), and shape every $T_o$ steps (often $T_o{=}1$).

\subsection{Sensor Response}
\label{subsec:sensor_response}

\paragraph{Response structure.}
$\mathsf{R}_y^n=\{z_{0:T_y}^{(y)},s_{0:T_y}^{(y)},\Delta\mathcal{E}_n\}$, where $\Delta\mathcal{E}_n$ contains newly detected obstacles
or refinements of obstacle parameters.

\paragraph{Active-set update.}
After $\mathcal{E}_n \leftarrow \mathcal{E}_n\cup\Delta\mathcal{E}_n$, we compute $\mathcal{C}_n$ using ~\cref{eq:impl:active_set}.

\paragraph{Local Hamiltonian.}
As in ~\cref{eq:module_template_hamiltonian}, the sensor module uses
\[
H_y(q_y,p_y;\mathcal{E}_n,\eta_n)
=\tfrac12 p_y^\top M_y^{-1}p_y+\mathcal{R}_y(q_y;\mathcal{E}_n,\eta_n),
\quad
\mathcal{R}_y=E_{\mathrm{sensor}}+\sum_{i\in\mathcal{C}_n}\alpha_{n,i}\,b(d_i;\hat d).
\]

\subsection{Frame Response}
\label{subsec:frame_response}

\paragraph{Response structure.}
$\mathsf{R}_f^n=\{z_{0:T_f}^{(f)},s_{0:T_f}^{(f)},\{v_\ell^{(f)}\}_{\ell=0}^{T_f}\}$, where $v_\ell^{(f)}=M_f^{-1}p_\ell^{(f)}$.

\paragraph{Local Hamiltonian with dissipation/port.}
The frame module follows ~\cref{eq:module_score_field} with
\[
H_f(q_f,p_f;\mathcal{E}_n,\eta_n)
=\tfrac12 p_f^\top M_f^{-1}p_f+\mathcal{R}_f(q_f;\mathcal{E}_n,\eta_n),
\quad
\mathcal{R}_f=\beta_n E_{\mathrm{goal}}+\sum_{i\in\mathcal{C}_n}\alpha_{n,i}\,b(d_i;\hat d),
\]
and $(\Gamma_f,G_f)$ instantiated by ~\cref{eq:impl:block_selectors}.

\subsection{Shape Response}
\label{subsec:shape_response}

\paragraph{Response structure.}
$\mathsf{R}_o^n=\{z_{0:T_o}^{(o)},s_{0:T_o}^{(o)},\{\min_i d_i(q_\ell^{(o)};\mathcal{E}_n)\}_{\ell=0}^{T_o}\}$.

\paragraph{Local Hamiltonian.}
\[
H_o(q_o,p_o;\mathcal{E}_n,\eta_n)
=\tfrac12 p_o^\top M_o^{-1}p_o+\mathcal{R}_o(q_o;\mathcal{E}_n,\eta_n),
\quad
\mathcal{R}_o=\lambda_n E_{\mathrm{obj}}+\sum_{i\in\mathcal{C}_n}\alpha_{n,i}\,b(d_i;\hat d),
\]
with no dissipation/ports on the shape subsystem (i.e., $\Gamma_o=0$, $u_o=0$).

\subsection{Hamiltonian Assembly and Discrete Integration}
\label{subsec:symplectic_integration}

\paragraph{Energy terms.}
We use the surrogate energy and Hamiltonian in ~\cref{eq:impl:surrogate_energy,eq:impl:surrogate_hamiltonian}.

\paragraph{One-step update.}
We implement the symplectic Euler step in ~\cref{eq:impl:ph_euler_p,eq:impl:ph_euler_q}, using gradients from
~\cref{eq:impl:gradients} and frame-only selectors from ~\cref{eq:impl:block_selectors}.

\subsection{Observable Extraction}
\label{subsec:observable_extraction}

We compute observables using ~\cref{eq:impl:observables,eq:impl:observable_vector} and compare against the target
in ~\cref{eq:impl:observable_target}.

\subsection{Parameter Adaptation via Observable Feedback}
\label{subsec:parameter_adaptation}

We form the active parameter vector via ~\cref{eq:impl:active_param_vector}, update $J_n$ via ~\cref{eq:impl:secant_J},
and solve the projected Tikhonov update via ~\cref{eq:impl:tikhonov_update}. Inactive barrier weights $\alpha_{n,j}$ for
$j\notin\mathcal{I}_n$ remain unchanged.

\subsection{Port Correction from Residual}
\label{subsec:port_correction}

We compute the residual via ~\cref{eq:impl:residual}, update $P_n$ via ~\cref{eq:impl:secant_P}, and solve the port LS with
clipping via ~\cref{eq:impl:port_update}. A robust fallback is to set $P_n=\mathrm{diag}(0,0,\kappa_v)$ (ports primarily regulate speed).

\subsection{Output Format}
\label{subsec:output_format}

\paragraph{Trajectory and parameters.}
Return $\{z_n\}_{n=0}^{N}$ and the history $\{(\eta_n,\mu_n,u_{f,n})\}_{n=0}^{N}$ where $N$ is the final step. This enables post-hoc
inspection of how $(\beta,\lambda,\alpha_i)$ and $\mu$ evolve with the active set in ~\cref{eq:impl:active_set}, and when nonconservative
port forcing is invoked via ~\cref{eq:impl:port_update}.

\subsection{Additional hyperparameters in Table~\ref{tab:hyperparams}}
\label{subsec:hyperparam}

As a complement to Algorithm~\ref{algo:online_grl_snam}, we specify additional fixed parameters and training hyperparameters below.

\paragraph{Surrogate training (offline).}
The surrogate coefficient network $g_\xi$ is trained on short rollouts of horizon $H \in \{2,\ldots,6\}$ steps with loss
\[
\mathcal{L} = w_{\mathrm{traj}}\|o_H - o_H^*\|^2 + w_{\mathrm{vel}}\|v_H - v_H^*\|^2 + w_{\mathrm{fric}}\|\gamma - \gamma_o\|^2 + w_{\mathrm{multi}}\mathcal{L}_{\mathrm{ms}},
\]
where $(w_{\mathrm{traj}}, w_{\mathrm{vel}}, w_{\mathrm{fric}}, w_{\mathrm{multi}}) = (1.0, 1.0, 0.1, 0.5)$ and $\mathcal{L}_{\mathrm{ms}}$ is a multi-start robustness penalty sampling perturbed initial conditions near obstacles. We use learning rate $\eta \in [10^{-4}, 3 \times 10^{-4}]$ with Adam and gradient clipping at norm $5.0$.

\paragraph{IPC barrier.}
The log-barrier $b(d; \hat{d}) = -(d - \hat{d})^2 \log(d/\hat{d})$ activates at distance $\hat{d} \in [0.28, 1.0]$. Barrier values and gradients are clamped to $[0, 200]$ and $[-200, 200]$ respectively for numerical stability.

\paragraph{Stage planning.}
Stages have dimensions $(W_{\mathrm{stage}}, H_{\mathrm{stage}}) = (2.6, 2.0)$ with overlap ratio $\rho_{\mathrm{overlap}} = 0.3$. Obstacles are inflated by $r_{\mathrm{inflate}} \in [0.05, 0.4]$ to define the collision-free tube for widest-path exit selection. The margin factor for radius-aware integration is $0.5$.

\paragraph{Hyperelastic ring experiment.}
\emph{Dynamics:} The ring uses mass matrix $M = \mathrm{blkdiag}(M_o I_2, I, M_s)$ with $(M_o, I, M_s) = (1.5, 0.6, 1.0)$ and nominal damping $(\gamma_o, \gamma_\theta, \gamma_s) = (4.0, 1.6, 2.0)$. Area resistance uses bulk modulus $k_{\mathrm{bulk}} = 1.5$. 
\emph{Geometry:} The periodic cubic B-spline boundary has $n_{\mathrm{ctrl}} = 20$ control points and $K = 240$ curve samples, with nominal robot radius $r_{\mathrm{robot}} \in [0.30, 0.50]$.
\emph{Integration:} Symplectic Euler with timestep $\tau = 0.03\,\mathrm{s}$.

\paragraph{Dungeon navigation experiment.}
\emph{Dataset generation:} Maximum velocity $v_{\max} \in [2.5, 3.0]\,\mathrm{px/s}$, goal attraction gain $k_{\mathrm{goal}} \in [20, 25]$, barrier repulsion gain $k_{\mathrm{barrier}} \in [0, 1]$ (often $0$ when SDF provides implicit avoidance), and barrier activation $\hat{d} = 1.0\,\mathrm{px}$.
\emph{Integration:} Explicit Euler with timestep $\tau = 0.3\,\mathrm{s}$ and $4$ substeps for collision resolution.
\emph{Environment:} Obstacle distances are queried from a precomputed signed distance field $\mathrm{SDF}(x,y)$ via Euclidean distance transform. Stages are square windows of side $s_{\mathrm{stage}} = 60\,\mathrm{px}$, with waypoints inflated to maintain wall clearance $d_{\mathrm{wp}} \geq 1.5$--$4.0\,\mathrm{px}$.

\paragraph{Test-time adaptation.}
When enabled, online finetuning uses proximal regularization $\lambda_{\mathrm{prox}} = 10^{-3}$ anchored to the pretrained checkpoint, updating only the final projection heads. The secant-based controller selects $k_\alpha \in \{1, 2\}$ nearest obstacles for local $\alpha$-adjustment with per-head learning rates $(\kappa_\beta, \kappa_\gamma, \kappa_\alpha) = (0.25, 0.05, 0.4)$.

